\section*{Chapitre \ref{ch:qm:stochastic-differential-equations} --- Équations différentielles stochastiques}

\begin{solution}{exo:qm:stochastic-differential-equations:1}
$\ln 2/2 = 0.35$ ans, soit 87 jours de bourse.
\end{solution}

\begin{solution}{exo:qm:stochastic-differential-equations:2}
$\P(\ln S_1 < \ln S_0) = \Phi(-(0.05 - 0.02)/0.20) = \Phi(-0.15) = 44.0\%$~: la médiane croît au taux $\mu -
\sigma^2/2 = 3\%$.
\end{solution}

\begin{solution}{exo:qm:stochastic-differential-equations:3}
$\mathcal Lf = \kappa(\bar x - x)f' + \tfrac12\sigma^2f^{\prime\prime}$~; pour $f(x) = x$, $\mathcal Lf = \kappa(\bar x - x)$, donc
$\frac{d}{dt}\E[X_t] = \kappa(\bar x - \E[X_t])$ et $\E[X_t] = \bar x + (X_0 - \bar x)e^{-\kappa t}$.
\end{solution}

\begin{solution}{exo:qm:stochastic-differential-equations:4}
$\kappa(\bar x - y)p = \tfrac12\sigma^2p'$ donne une densité gaussienne de variance $\sigma^2/2\kappa = 0.16/4 = 0.04$.
\end{solution}

\begin{solution}{exo:qm:stochastic-differential-equations:5}
La formule de Dynkin avec $f(v) = v$ donne $\frac{d}{dt}\E[v_t] = \kappa(\bar v - \E[v_t])$, donc $\E[v_{0.5}] = 0.04 + 0.05e^{-1} =
0.0584$.
\end{solution}

\begin{solution}{exo:qm:stochastic-differential-equations:6}
$B = (1 - e^{-2.5})/0.5 = 1.8358$~; $A = (0.04 - 0.0002)(1.8358 - 5) - 0.0001 \times 1.8358^2/2 = -0.1261$~;
$e^{A - 0.03B} = e^{-0.1812} = 0.8343$.
\end{solution}

\begin{solution}{exo:qm:stochastic-differential-equations:7}
24.5~\% pour un ratio de Feller égal à un~: la condition tient le processus exact à l'écart de zéro, et non le
pas gaussien d'Euler. Dans le tableau, $\eta = 0.25$ (ratio 2.56) donne 0.08~\%, et $\eta = 0.3$ (ratio 1.78)
déjà 1.6~\%.
\end{solution}

\begin{solution}{exo:qm:stochastic-differential-equations:8}
Tronquer à zéro supprime les NaN mais non l'erreur~: chaque pas qui aurait franchi zéro est remplacé par
zéro, qui a alors une diffusion nulle, si bien que le schéma ajoute de la masse en zéro et biaise la loi près de la
frontière. C'est toujours un schéma d'Euler d'ordre un, et non la transition exacte~; vérifier ses moments
par rapport à $\E[v_t] = \bar v + (v_0 - \bar v)e^{-\kappa t}$ et à la \omterm{def:qm:stochastic-differential-equations:kolmogorov}{loi stationnaire} gamma, ou utiliser le schéma exact ou
le schéma à troncature complète.
\end{solution}

\begin{solution}{pb:qm:stochastic-differential-equations:1}
\textbf{1.} $2 \times 2 \times 0.04/0.36 = 0.44$~: inférieur à un, zéro est atteignable.
\textbf{2.} $\mathrm{Gamma}(0.44, 0.09)$~: moyenne 0.04, écart-type $\sqrt{0.04 \times 0.36/4} = 0.06$.
\textbf{3.} 27.9~\%.
\textbf{4.} $\eta = \sqrt{2 \times 2 \times 0.04} = 0.4$.
\textbf{5.} $\ln 2/2 = 0.35$ ans.
\textbf{6.} $\Phi(-(0.004 + 2 \times 0.036/252)/(0.6\sqrt{0.004/252})) = 3.6\%$.
\textbf{7.} 75.2~\% des 40\,000 trajectoires.
\textbf{8.} 75.3~\% avec quatre pas par jour, 74.6~\% avec seize, 75.2~\% avec des pas hebdomadaires.
\textbf{9.} Le processus exact s'approche arbitrairement de zéro sur ces trajectoires~; à partir d'un petit $v$, un
pas gaussien, quelle que soit sa durée, a une probabilité fixe de franchir zéro, et davantage de pas près de zéro signifie davantage de
chances de franchissement.
\textbf{10.} Il se propage à travers chaque somme et chaque moyenne jusqu'à la valeur du portefeuille, et une comparaison
avec un NaN est fausse, si bien que les contrôles de limite peuvent passer sans rien signaler.
\textbf{11.} Minimum 0 (jamais négatif)~; moyenne annuelle de 0.0394, à comparer avec $\bar v = 0.04$.
\textbf{12.} Moyenne 0.0393~; nulle sur 5.2~\% des pas.
\textbf{13.} Le schéma exact n'a aucun biais de discrétisation~; la troncature complète en a un, et son biais
s'annule lorsque $\Delta t \to 0$.
\textbf{14.} La réflexion transforme chaque dépassement en une valeur positive de même taille, ce qui repousse la masse
vers le haut et biaise la moyenne à la hausse près de zéro.
\textbf{15.} Le tirage exact est lent et les variantes d'Euler exigent de petits pas~; le
schéma quadratique-exponentiel (livre~5, chapitre~23) est précis pour des pas journaliers ou plus grands.
\textbf{16.} Non~: le sourire de volatilité du marché exige une forte volatilité de la volatilité, et imposer la condition
fausserait l'ajustement pour sauver un schéma numérique~; il faut corriger le schéma.
\textbf{17.} Aucun NaN ni variance négative, quelle que soit la graine~; $\E[v_t]$ et $\Var(v_t)$ comparées à leurs
formes fermées~; la \omterm{def:qm:stochastic-differential-equations:kolmogorov}{loi stationnaire} comparée à la densité gamma~; la convergence lorsque le pas est raffiné.
\textbf{18.} La cause première (Euler simple avec un ratio de Feller de 0.44), l'ampleur de la défaillance (75~\% des
trajectoires), le correctif (schéma exact ou à troncature complète) et sa validation, ainsi que les tests ajoutés.
\textbf{19.} Résultat nommé~: \emph{le NaN de la passe de nuit}~: avec un ratio de Feller de 0.44, 75~\% des
trajectoires d'Euler simple au pas quotidien produisent une variance négative en moins d'un an, quel que soit le pas~; le
schéma exact et le schéma à troncature complète n'en produisent aucune.
\textbf{20.} Elle tient le processus exact à l'écart de zéro~; elle n'empêche pas un pas gaussien de
franchir zéro.
\end{solution}

\begin{solution}{iq:qm:stochastic-differential-equations:1}
$S_T = S_0\exp((\mu - \tfrac12\sigma^2)T + \sigma W_T)$~; la médiane vaut $S_0e^{(\mu - \sigma^2/2)T}$, inférieure à la moyenne
$S_0e^{\mu T}$.
\iqlookfor{Itô sur $\ln S$ et la loi lognormale.}
\end{solution}

\begin{solution}{iq:qm:stochastic-differential-equations:2}
$\kappa = \ln 2/10 = 0.069$ par jour~; au bout de 30 jours, soit trois \omterm{def:qm:stochastic-differential-equations:ou}{demi-vies}, il reste $2^{-3} = 12.5\%$.
\iqlookfor{$e^{-\kappa t}$ et \omterm{def:qm:stochastic-differential-equations:ou}{demi-vies}.}
\end{solution}

\begin{solution}{iq:qm:stochastic-differential-equations:3}
$2\kappa\bar v \ge \eta^2$~: alors la dérive près de zéro l'emporte sur la diffusion et zéro n'est jamais atteint~;
sinon le processus touche zéro et se réfléchit. La loi reste positive ou nulle dans les deux cas, mais un
pas d'Euler partant du voisinage de zéro peut devenir négatif, et davantage lorsque la condition échoue.
\iqlookfor{la classification de la frontière et sa conséquence numérique.}
\end{solution}

\begin{solution}{iq:qm:stochastic-differential-equations:4}
$u(t,x) = \E[e^{-\int_t^Tr}g(X_T) \mid X_t = x]$ est solution de $\partial_tu + \mathcal Lu - ru = 0$, $u(T) = g$. Pour
$dS = \mu S\,dt + \sigma S\,dW$ et $r$ constant~: $\partial_tu + \mu S\partial_Su + \tfrac12\sigma^2S^2\partial_{SS}u - ru = 0$.
\iqlookfor{générateur, actualisation et condition terminale.}
\end{solution}

\begin{solution}{iq:qm:stochastic-differential-equations:5}
Rendre la défaillance reproductible (consigner la graine et l'indice de trajectoire), trouver la première valeur non finie
et le pas qui l'a produite, puis lire le schéma dans cet état. Pour la rendre impossible~: schémas exacts ou
schémas respectant la frontière, assertions sur l'espace d'état dans le noyau, un test qui exécute de nombreuses
graines et échoue sur tout NaN, et une agrégation qui refuse les entrées non finies au lieu de
les propager.
\iqlookfor{reproductibilité, cause première et invariants imposés dans le code.}
\end{solution}

\begin{solution}{iq:qm:stochastic-differential-equations:6}
Une \omterm{def:qm:stochastic-differential-equations:sde}{solution forte} est construite sur un \omterm{def:qm:brownian-motion:bm}{mouvement brownien} donné~; une \omterm{def:qm:stochastic-differential-equations:sde}{solution faible} demande seulement un espace de probabilité
portant un processus et un \omterm{def:qm:brownian-motion:bm}{mouvement brownien} qui vérifient l'équation. L'équation de Tanaka $dX =
\operatorname{sign}(X)\,dW$, $X_0 = 0$, admet des \omterm{def:qm:stochastic-differential-equations:sde}{solutions faibles} (tout \omterm{def:qm:brownian-motion:bm}{mouvement brownien} $X$ avec $W = \int\operatorname{sign}(X)\,dX$)
mais aucune \omterm{def:qm:stochastic-differential-equations:sde}{solution forte}, car $X$ ne peut pas être reconstruit à partir de $W$.
\iqlookfor{les définitions et l'exemple de Tanaka.}
\end{solution}
